% VITAL -- draft abstract package for the professor (structure, gap, research question)
% Compile: pdflatex VITAL_abstract_draft.tex
\documentclass[11pt]{article}
\usepackage[margin=1in]{geometry}
\usepackage{booktabs,array,amssymb,pifont,enumitem,xcolor,hyperref}
\usepackage[numbers]{natbib}
\newcommand{\yes}{\ding{51}}            % full support
\newcommand{\prt}{$\circ$}             % partial / via add-on
\newcommand{\no}{--}                    % not addressed
\newcommand{\todo}[1]{\textcolor{red}{[#1]}}
\setlist{itemsep=2pt,topsep=3pt}

\title{\textbf{VITAL: A Verified, Frame-Explicit Flight Simulation Framework\\
that Turns Aircraft Models into Traceable Flying-Qualities Verdicts}\\[4pt]
\large Draft abstract and paper outline}
\author{Krish Bhatt \\ Virginia Tech}
\date{\today}

\begin{document}
\maketitle

%=====================================================================
\section*{Draft abstract (about 250 words)}
Flying-qualities problems are usually discovered late, after a configuration
is frozen, because the tools that model an aircraft, the tools that simulate
it, and the documents that judge it are disconnected. Open flight-dynamics
codes simulate well but do not issue verdicts against handling-qualities
requirements; conceptual-design environments predict stability but are
rarely verified against independent reference data; and clearance methods
quantify robustness but are tied to one aircraft and one control law.
Certification by analysis needs all three at once: a credible model, a
verified simulation, and a rule-based verdict whose margin and uncertainty
are stated.

This work presents VITAL, an aircraft-agnostic framework that links these
steps: concept, model, trim and linearization, rule verdict, and design
feedback. Aircraft are defined as data:
\begin{itemize}
  \item an explicit tree of reference axes;
  \item load-producing components;
  \item mass properties, read from standard DAVE-ML (ANSI/AIAA S-119) files.
\end{itemize}
Every frame transfer and moment-reference change is therefore performed by
one verified mechanism.

Each capability is accepted only through gated, test-first evidence:
\begin{itemize}
  \item published NASA check cases;
  \item independent toolbox results;
  \item analytic solutions;
  \item deliberate code sabotages that the tests must detect.
\end{itemize}

The framework is demonstrated on the NASA NESC F-16 model. Early results:
\begin{itemize}
  \item the published trim is reproduced to within $0.0004^\circ$;
  \item a 10-second Aerospace Blockset 6-DOF hold stays at trim to
    $10^{-14}$, and agrees with the framework's own equations of motion
    to $10^{-12}$;
  \item MIL-F-8785C Class IV levels are assessed with stated margins.
\end{itemize}
The paper then shows how uncertainty-aware verdicts and sensitivity-based
control-law changes restore compliance. It outlines the extension to
distributed electric propulsion (NASA X-57) and to FAA Part 23 criteria.
\todo{update numbers at submission}

%=====================================================================
\section{Main gap in the literature}
\begin{itemize}
  \item \textbf{Simulation and judgement are separate.} Mature simulators
    (JSBSim~\citep{berndt2004}, MATLAB Aerospace Blockset) integrate the
    equations of motion but stop at time histories. Turning those into a
    \emph{verdict} against MIL-F-8785C or Part 23 criteria is left to
    project-specific scripts.
  \item \textbf{Design-stage stability tools are rarely verified against
    independent data.}
    \begin{itemize}
      \item CEASIOM/SDSA~\citep{rizzi2011simsac} and total-envelope
        conceptual S\&C methods~\citep{swann2015} predict modes and levels
        early.
      \item SUAVE~\citep{lukaczyk2015} supports unconventional
        configurations.
      \item None of them reports reproduction of public multi-simulation
        check cases such as NASA/TM-2015-218675~\citep{nesc2015}.
    \end{itemize}
  \item \textbf{Robust clearance is aircraft- and law-specific.} GARTEUR
    clearance methods~\citep{fielding2002} quantify worst cases, but they
    are not embedded in a reusable model-to-verdict pipeline. Their
    ``pass'' semantics under uncertainty are not standardized.
  \item \textbf{Credibility is the bottleneck for certification by
    analysis.} NASA's CbA guide~\citep{mauery2021} and recent credibility
    studies~\citep{favaro2025} identify verified models, quantified
    uncertainty, and traceable evidence as the missing pieces.
  \item \textbf{Frame and reference-point errors are silent.} Moment-reference
    shifts, canted thrust lines, station/body conventions and inertia
    products are typically hand-coded per aircraft. For distributed-propulsion
    aircraft with many canted propulsors, this does not scale and is not
    auditable.
\end{itemize}

\begin{table}[h]
\centering\small
\caption{Capabilities of representative tools as described in their
published sources (\yes\ addressed, \prt\ partial or via add-on,
\no\ not addressed). \todo{confirm each cell during the literature review}}
\label{tab:gap}
\begin{tabular}{@{}>{\raggedright}p{3.1cm}*{8}{>{\centering\arraybackslash}p{1.25cm}}@{}}
\toprule
Tool / approach & Nonlinear 6-DOF sim & Public check-case verif. & Explicit frames \& ref.\ points & Trim / lin.\ / modes & Rule-based FQ verdict & Uncert.-aware verdict & Design feedback & Any config.\ (e.g.\ DEP) \\
\midrule
JSBSim~\citep{berndt2004}             & \yes & \prt & \prt & \prt & \no   & \no   & \no   & \yes \\
Aerospace Blockset                    & \yes & \no   & \prt & \prt & \no   & \no   & \no   & \prt \\
SUAVE~\citep{lukaczyk2015}            & \no  & \no   & \prt & \prt & \no   & \no   & \yes  & \yes \\
CEASIOM / SDSA~\citep{rizzi2011simsac}& \yes & \no   & \prt & \yes  & \yes  & \no   & \prt & \prt \\
Total-envelope S\&C~\citep{swann2015} & \no  & \no   & \no   & \yes  & \prt & \no   & \yes  & \prt \\
GARTEUR clearance~\citep{fielding2002}& \yes & \no   & \no   & \yes  & \yes  & \yes  & \no   & \no \\
NESC check cases~\citep{nesc2015}     & \yes & \yes  & \no   & \no   & \no   & \no   & \no   & \no \\
\midrule
\textbf{VITAL (this work)}            & \yes & \yes  & \yes  & \yes  & \yes  & \yes  & \yes  & \prt\,$\rightarrow$\,\yes \\
\bottomrule
\end{tabular}
\end{table}

\paragraph{Gap statement (one sentence).}
No openly documented framework carries an arbitrary aircraft definition
through a simulation that has been verified against independent public
reference data, all the way to a rule-based flying-qualities verdict. Such a
verdict would state its margin and uncertainty, and would identify the design
change that restores compliance.

%=====================================================================
\section{Research question}
\textbf{Main question.} Can an aircraft-agnostic simulation framework, verified
against independent public reference data, produce \emph{credible and
traceable} flying-qualities verdicts, with explicit margins and uncertainty,
early enough in design to guide configuration and control-law changes?

\textbf{Sub-questions.}
\begin{enumerate}[label=RQ\arabic*.]
  \item \textbf{Representation.} Can an explicit tree of axes, components and
    mass properties represent conventional and distributed-propulsion aircraft
    with no aircraft-specific plant code? Can it do so without loss of
    accuracy or speed against a hand-written model?
  \item \textbf{Verification.} What evidence is sufficient for each
    capability to count as credible? The candidates are public check cases,
    independent tools, analytic solutions and mutation (sabotage) testing.
  \item \textbf{Verdict.} How should a ``pass'' be defined when model
    parameters are uncertain? Specifically, how should worst-case margin and
    probability of meeting a level be combined so that a robust pass is never
    weaker than a nominal pass?
  \item \textbf{Feedback.} Can sensitivities of the verdict margin identify
    the design or control-law change that restores compliance, confirmed by
    re-assessment?
\end{enumerate}

%=====================================================================
\section{Proposed paper structure}
\begin{enumerate}[label=\arabic*.,leftmargin=*]
\item \textbf{Introduction.}
  \begin{itemize}
    \item Why flying-qualities issues surface late.
    \item The push toward certification by analysis.
    \item New configurations (distributed electric propulsion, eVTOL).
    \item Contributions, and the research question.
  \end{itemize}

\item \textbf{Related work.}
  \begin{enumerate}[label*=\arabic*.]
    \item \emph{Flight-dynamics simulators}: JSBSim, Aerospace Blockset;
      what they verify and what they leave to the user.
    \item \emph{Conceptual-design stability and control}: CEASIOM/SDSA,
      SUAVE, total-envelope methods.
    \item \emph{Flight-control clearance and robustness}: GARTEUR AG11,
      $\mu$-analysis, optimization-based clearance.
    \item \emph{Simulation credibility and certification by analysis}:
      NASA CbA guide, NASA-STD-7009 \todo{cite},
      UQ-based credibility.
    \item \emph{Gap}: Table~\ref{tab:gap} and the gap statement.
  \end{enumerate}

\item \textbf{The VITAL framework.}
  \begin{enumerate}[label*=\arabic*.]
    \item \emph{Frame-explicit object core.}
      \begin{itemize}
        \item Axis tree, points, vectors, loads (wrenches), inertia.
        \item One verified mechanism for every frame and moment-reference
          transfer.
      \end{itemize}
    \item \emph{Standards-based data layer.}
      \begin{itemize}
        \item DAVE-ML import with units and provenance.
        \item Hash-protected source data.
        \item Embedded check data reproduced.
      \end{itemize}
    \item \emph{Aircraft as data.}
      \begin{itemize}
        \item Components (aero database, propulsors) on axes, plus mass
          properties.
        \item A compiled fast path proven identical to the object path.
      \end{itemize}
    \item \emph{Analyses.}
      \begin{itemize}
        \item Typed trim with status codes and stall/data-envelope bounds.
        \item Linearization with step-plateau check.
        \item Mode identification.
      \end{itemize}
    \item \emph{Rule engine.} Cited, machine-readable criteria records
      (MIL-F-8785C now, FAA Part 23 later); levels and normalized margins
      over the flight envelope.
    \item \emph{Uncertainty and verdict logic.}
      \begin{itemize}
        \item Bound-worst margin over a joint-coverage box.
        \item Monte Carlo at the critical condition.
        \item A ``robust'' verdict defined so that it is never weaker than
          ``pass''.
      \end{itemize}
    \item \emph{Simulink and real time.}
      \begin{itemize}
        \item Generated 6-DOF models driven by the same aircraft object.
        \item Visualization and pilot-in-the-loop.
      \end{itemize}
  \end{enumerate}

\item \textbf{Verification methodology.}
  \begin{itemize}
    \item Evidence classes: published, independent, analytic, regression.
    \item Red-green gated development.
    \item A failure catalogue mapped to tests.
    \item Sabotage (mutation) testing: every gate must detect deliberate
      code faults.
    \item NESC check cases as the external reference.
  \end{itemize}

\item \textbf{Case study: NASA F-16.}
  \begin{enumerate}[label*=\arabic*.]
    \item Data provenance: NASA TP-1538 $\rightarrow$ Stevens \& Lewis
      $\rightarrow$ NESC DAVE-ML.
    \item Trim reproduction and the 6-DOF trim hold in Simulink.
    \item NESC dynamic check cases within the multi-simulation band.
    \item MIL-F-8785C Class IV levels, critical conditions, pre-registered
      degradations that must fail.
    \item Design feedback: a pitch SAS and yaw damper restore the target
      level, confirmed by re-assessment.
    \item Uncertainty: robust verdicts and their sensitivity to assumed
      widths.
  \end{enumerate}

\item \textbf{Toward distributed electric propulsion (planned).}
  \begin{itemize}
    \item A twin-engine and engine-out known-answer aircraft.
    \item NASA X-57 Maxwell: canted propulsors, gyroscopic and propeller-wing
      coupling.
    \item Part 23 criteria.
  \end{itemize}

\item \textbf{Discussion and limitations.}
  \begin{itemize}
    \item Judgement-based uncertainty widths.
    \item Linear criteria applied to nonlinear aircraft.
    \item The scope of MIL-F-8785C versus modern standards.
    \item The cost of verification.
  \end{itemize}

\item \textbf{Conclusions and future work.}
\end{enumerate}

%=====================================================================
\section*{Expected contributions}
\begin{enumerate}[label=C\arabic*.]
  \item An aircraft-agnostic, frame-explicit model representation, verified to
    reproduce a hand-written reference model exactly.
  \item A gated, evidence-tagged verification method, including sabotage
    testing, demonstrated against public NASA check cases.
  \item A rule-based flying-qualities verdict with explicit margins, and a
    consistent definition of a robust pass under uncertainty.
  \item A demonstrated loop from verdict to design change on the F-16, with a
    path to distributed-propulsion aircraft.
\end{enumerate}

%=====================================================================
\begin{thebibliography}{99}\small
\bibitem[Berndt(2004)]{berndt2004} J.~S. Berndt, ``JSBSim: An Open Source
  Flight Dynamics Model in C++,'' AIAA Modeling and Simulation Technologies
  Conf., AIAA 2004-4923, 2004. doi:10.2514/6.2004-4923.
\bibitem[Lukaczyk et~al.(2015)]{lukaczyk2015} T.~W. Lukaczyk et~al., ``SUAVE:
  An Open-Source Environment for Multi-Fidelity Conceptual Vehicle Design,''
  16th AIAA/ISSMO MA\&O Conf., AIAA 2015-3087, 2015.
\bibitem[Rizzi et~al.(2011)]{rizzi2011simsac} A.~Rizzi et~al., ``Modeling and
  Simulating Aircraft Stability and Control---The SimSAC Project,'' Progress
  in Aerospace Sciences, 2011. \todo{verify details; add
  Goetzendorf-Grabowski SDSA paper}
\bibitem[Swann and Takahashi(2015)]{swann2015} J.~Swann and T.~T. Takahashi,
  ``A Total Flight Envelope Approach to Conceptual Design Stability \&
  Control,'' AIAA 2015-3377, 2015.
\bibitem[Fielding et~al.(2002)]{fielding2002} C.~Fielding, A.~Varga,
  S.~Bennani, M.~Selier (eds.), \emph{Advanced Techniques for Clearance of
  Flight Control Laws}, LNCIS 283, Springer, 2002.
\bibitem[Murri et~al.(2015)]{nesc2015} D.~G. Murri, E.~B. Jackson,
  R.~O. Shelton et~al., ``Check-Cases for Verification of
  6-Degree-of-Freedom Flight Vehicle Simulations,'' NASA/TM-2015-218675,
  Vols.~I--II, 2015.
\bibitem[Mauery and Alonso(2021)]{mauery2021} T.~Mauery, J.~Alonso et~al.,
  ``A Guide for Aircraft Certification by Analysis,'' NASA CR, NTRS
  20210015404, 2021.
\bibitem[Favaro et~al.(2025)]{favaro2025} F.~Favaro, D.~Rylko, G.~Quaranta,
  ``Certification by Simulation of an eVTOL: Credibility Analysis of
  Simulation Model Using Uncertainty Quantification,'' European Rotorcraft
  Forum, 2025.
\bibitem[Nguyen et~al.(1979)]{nguyen1979} L.~T. Nguyen et~al., ``Simulator
  Study of Stall/Post-Stall Characteristics of a Fighter Airplane with
  Relaxed Longitudinal Static Stability,'' NASA TP-1538, 1979.
\bibitem[Stevens and Lewis(2003)]{stevens2003} B.~L. Stevens and
  F.~L. Lewis, \emph{Aircraft Control and Simulation}, 2nd ed., Wiley, 2003.
\bibitem[MIL-F-8785C(1980)]{mil8785c} \emph{Military Specification: Flying
  Qualities of Piloted Airplanes}, MIL-F-8785C, 1980.
\bibitem[AIAA(2011)]{s119} ANSI/AIAA S-119-2011, \emph{Flight Dynamic Model
  Exchange Standard} (DAVE-ML).
\end{thebibliography}
\end{document}
